% Pattern: number/argument. A wrapper \newenvironment{restate}[1] with
% \def\thetheorem{#1} prints the number of \ref{thm:main} for its argument,
% and the \addtocounter{theorem}{-1} of its end code takes back the step, so
% the theorem after it is 2.2. Seen in the gold pilot (set X).
% want theorem-1 statement/theorem 1.1
% want theorem-2 statement/theorem 2.1
% want restate-1 statement/theorem 1.1
% want theorem-3 statement/theorem 2.2
\documentclass{article}
\usepackage{amsthm}
\newtheorem{theorem}{Theorem}[section]
\newenvironment{restate}[1]
    {\def\thetheorem{#1}\begin{theorem}}
    {\end{theorem}\addtocounter{theorem}{-1}}
\begin{document}
\section{One}
\begin{theorem}\label{thm:main} Every finite space is compact. \end{theorem}
\section{Two}
\begin{theorem} Every point is closed. \end{theorem}
\begin{restate}{\ref{thm:main}} Every finite space is compact. \end{restate}
\begin{theorem} Every closed set is compact. \end{theorem}
\end{document}
